Los datos de entrenamiento son una parte crucial del desarrollo del modelo. El \textit{dataset} , concepto ya definido en la sección correspondiente, corresponde a la información de la que parte el sistema de aprendizaje para inferir información nueva, por lo que su elección es crucial para que el modelo desempeñe correctamente la tare especificada.

Para la tarea de separación de fuentes los datos de entrenamiento elegidos han sido tomados de la base de datos de \textit{MUSDB18} \cite{musdb18}. El conjunto ya se ha mencionado antes en esta memoria, y se compone de 150 canciones completas de diferentes géneros junto con sus fuentes separadas: \textit{drums}, \textit{bass}, \textit{vocals} y \textit{others}.
El conjunto se separa en dos carpetas, \textit{train}, de 100 canciones y \textit{test}, de 50 canciones respectivamente. La documentación indica que los enfoques supervisados deben ser entrenados en el conjunto de entenamiento y probados en los conjuntos de test y de entrenamiento, y, además, que todas las pistas son estereofónicas y codificadas en 44.1kHz.